\uuid{UKy7}
\titre{DL usuels/Étude locale}
\auteur{hanani ; mzari}
\organisation{exo7}
\licence{CC BY-NC-SA 4.0}
\datecreate{2020-01-20}
\niveau{L1}
\module{Analyse}
\chapitre{Développement limité}
\sousChapitre{Applications}
\difficulte{3}
\theme{Applications des DL, DL usuels, Étude locale}
\exo7qcmid{lille-2.developpements-limites.043}

\contenu{
\texte{
Soit $a\in \Rr$ et $f$ une fonction. On dit que $f(x)$ est un $o\left((x-a)^n\right)$ au voisinage de $a$ et on écrit $f(x)=o((x-a)^n)$ si $\displaystyle \lim _{x\to a}\frac{f(x)}{(x-a)^n}=0$. On note par $DL_n(a)f(x)$ le développement limité de $f$ à l'ordre $n$ en $a$. On note aussi $DL_n(a^+)f(x)$ (resp. $DL_n(a^-)f(x)$) le développement limité de $f$ à l'ordre $n$ à droite en $a$ (resp. à gauche en $a$).
}

\begin{question}
\qclassification{125.03}
\qtags{motcle=DL usuels/Étude locale}
Soit $f$ telle que $\displaystyle f(x)=\frac{2}{\sin (x^2)}+\frac{2}{\ln (1-x^2)}$ si $x\neq 0$ et $f(0)=1$. La tangente au graphe de $f$ au point d'abscisse $0$ lorsqu'elle existe est notée $T_0$. Parmi les affirmations suivantes, cocher celles qui sont vraies :
\begin{answers}
\bad{$f$ n'est pas dérivable en $0$.}
\bad{$f$ n'admet pas de développement limité d'ordre $2$ en $0$.}
\good{$\displaystyle DL_2(0)f(x)=1+\frac{x^2}{2}+o(x^2)$.}
\good{$T_0$ est la droite d'équation $\displaystyle y=1$ et le graphe de $f$ est au dessus de $T_0$ au voisinage de $0$.}
\end{answers}
\vskip2mm
\begin{explanations}
D'abord, $\displaystyle f(x)=2\frac{\ln (1-x^2)+\sin (x^2)}{\sin (x^2)\ln (1-x^2)}$, et pour écrire le $DL_2(0)f(x)$, on utilise les $DL_6(0)\sin (x^2)$ et $DL_6(0)\ln(1-x^2)$. On trouve
$$\displaystyle DL_2(0)f(x)=1+\frac{x^2}{2}+o(x^2).$$
Ainsi $f$ est dérivable en $0$, $T_0$ est la droite d'équation $\displaystyle y=1$ et puisque $\displaystyle f(x)-y\simeq \frac{x^2}{2}\geq 0$ au voisinage de $0$ : le graphe de $f$ est au dessus de $T_0$ au voisinage de $0$.
\end{explanations}
\end{question}
}
